\documentclass[12pt,a4paper]{article}

\usepackage[margin=1in]{geometry}
\usepackage{setspace}
\usepackage{graphicx}
\usepackage{xcolor}
\usepackage{hyperref}
\usepackage{booktabs}
\usepackage{array}
\usepackage{float}
\usepackage{verbatim}

\onehalfspacing

\definecolor{boxgray}{rgb}{0.93,0.93,0.93}

% A simple shaded definition box
\newcommand{\defbox}[2]{%
  \vspace{0.2cm}
  \noindent\colorbox{boxgray}{\parbox{\dimexpr\linewidth-2\fboxsep}{%
    \textbf{#1:} #2
  }}
  \vspace{0.2cm}
}

% Placeholder box for screenshots to be inserted after completing the lab
\newcommand{\screenshotplaceholder}[2]{%
  \begin{figure}[H]
      \centering
      \fbox{\parbox[c][#1][c]{0.9\linewidth}{\centering #2}}
  \end{figure}
}

\begin{document}

% -------------------------------------------------------------
% Header
% -------------------------------------------------------------
\begin{center}
    {\Large \textbf{CSE312: Database Management Systems Lab}}\\[0.3cm]
    {\large \textbf{Lab 02: PostgreSQL DDL Statements}}
\end{center}

\vspace{0.5cm}

% -------------------------------------------------------------
% 1. Objective
% -------------------------------------------------------------
\section*{1. Objective}
The objectives of this experiment are:
\begin{itemize}
    \item To understand the purpose of Data Definition Language (DDL) in PostgreSQL.
    \item To create and manage databases, schemas, and tables using PostgreSQL commands.
    \item To define table columns with suitable data types and constraints.
    \item To modify existing database objects using \texttt{ALTER} statements.
    \item To remove database objects safely using \texttt{DROP} and \texttt{TRUNCATE} statements.
\end{itemize}

% -------------------------------------------------------------
% 2. Introduction
% -------------------------------------------------------------
\section*{2. Introduction}
Data Definition Language (DDL) is a part of SQL used to define, modify, and remove database
structures. In PostgreSQL, DDL statements are used to create databases, schemas, tables, constraints,
indexes, views, and other database objects.

This lab focuses on the most commonly used PostgreSQL DDL statements:
\begin{itemize}
    \item \texttt{CREATE} -- creates a new database object.
    \item \texttt{ALTER} -- modifies an existing database object.
    \item \texttt{DROP} -- permanently removes a database object.
    \item \texttt{TRUNCATE} -- removes all rows from a table while keeping the table structure.
\end{itemize}

\noindent\textbf{Required software:}
\begin{itemize}
    \item PostgreSQL Server
    \item pgAdmin or psql command-line tool
\end{itemize}

\screenshotplaceholder{5cm}{Screenshot Placeholder: PostgreSQL server or pgAdmin dashboard}

% -------------------------------------------------------------
% 3. Theory
% -------------------------------------------------------------
\section*{3. Theory}

\subsection*{3.1 Data Definition Language}
\defbox{DDL}{Data Definition Language is the subset of SQL used to define and manage the structure of a database. DDL commands affect database objects such as databases, schemas, tables, and constraints.}

\noindent Common DDL statements are summarized below:

\begin{table}[H]
\centering
\renewcommand{\arraystretch}{1.4}
\begin{tabular}{>{\bfseries}p{3cm} p{9cm}}
\toprule
Statement & Purpose \\
\midrule
CREATE & Creates a database object such as a database, schema, or table. \\
ALTER & Modifies the structure of an existing database object. \\
DROP & Deletes a database object permanently. \\
TRUNCATE & Removes all rows from a table without deleting the table itself. \\
\bottomrule
\end{tabular}
\caption{Common PostgreSQL DDL statements}
\label{tab:ddl_statements}
\end{table}

\subsection*{3.2 PostgreSQL Data Types}
Choosing appropriate data types is important because data types control what kind of values can be
stored in each column.

\begin{table}[H]
\centering
\renewcommand{\arraystretch}{1.4}
\begin{tabular}{>{\bfseries}p{4cm} p{8cm}}
\toprule
Data Type & Description \\
\midrule
INTEGER & Stores whole numbers. \\
SERIAL & Auto-incrementing integer, commonly used for primary keys. \\
VARCHAR(n) & Stores variable-length character strings with a maximum length. \\
TEXT & Stores variable-length text without a fixed maximum length. \\
NUMERIC(p, s) & Stores exact decimal values with precision and scale. \\
DATE & Stores calendar dates. \\
TIMESTAMP & Stores date and time values. \\
BOOLEAN & Stores true or false values. \\
\bottomrule
\end{tabular}
\caption{Frequently used PostgreSQL data types}
\label{tab:data_types}
\end{table}

\subsection*{3.3 Constraints}
\defbox{Constraint}{A constraint is a rule applied to a table column or a group of columns to maintain data accuracy and integrity.}

\noindent Common constraints include:
\begin{itemize}
    \item \texttt{PRIMARY KEY}: uniquely identifies each row in a table.
    \item \texttt{FOREIGN KEY}: links a column to the primary key of another table.
    \item \texttt{NOT NULL}: prevents a column from storing null values.
    \item \texttt{UNIQUE}: ensures all values in a column are different.
    \item \texttt{CHECK}: ensures values satisfy a specified condition.
    \item \texttt{DEFAULT}: assigns a value automatically when no value is provided.
\end{itemize}

% -------------------------------------------------------------
% 4. Procedure
% -------------------------------------------------------------
\section*{4. Procedure}

\subsection*{4.1 Connecting to PostgreSQL}
Open pgAdmin or the psql command-line tool and connect to the PostgreSQL server using a valid
username and password.

\begin{verbatim}
psql -U postgres
\end{verbatim}

\screenshotplaceholder{5cm}{Screenshot Placeholder: Successful connection to PostgreSQL}

\subsection*{4.2 Creating a Database}
Use the \texttt{CREATE DATABASE} statement to create a new database for this lab.

\begin{verbatim}
CREATE DATABASE university_lab;
\end{verbatim}

After creating the database, connect to it.

\begin{verbatim}
\c university_lab
\end{verbatim}

\screenshotplaceholder{5cm}{Screenshot Placeholder: Database creation and connection output}

\subsection*{4.3 Creating a Schema}
A schema is a logical container for database objects. It helps organize tables and other objects.

\begin{verbatim}
CREATE SCHEMA academics;
\end{verbatim}

Set the search path so PostgreSQL uses this schema by default.

\begin{verbatim}
SET search_path TO academics;
\end{verbatim}

\subsection*{4.4 Creating Tables}
The \texttt{CREATE TABLE} statement defines a new table, its columns, data types, and constraints.

\begin{verbatim}
CREATE TABLE department (
    dept_id SERIAL PRIMARY KEY,
    dept_name VARCHAR(100) NOT NULL UNIQUE,
    location VARCHAR(100)
);
\end{verbatim}

\begin{verbatim}
CREATE TABLE student (
    student_id SERIAL PRIMARY KEY,
    first_name VARCHAR(50) NOT NULL,
    last_name VARCHAR(50) NOT NULL,
    email VARCHAR(100) UNIQUE,
    date_of_birth DATE,
    admission_year INTEGER CHECK (admission_year >= 2000),
    dept_id INTEGER REFERENCES department(dept_id)
);
\end{verbatim}

\begin{verbatim}
CREATE TABLE course (
    course_id SERIAL PRIMARY KEY,
    course_code VARCHAR(20) NOT NULL UNIQUE,
    course_title VARCHAR(150) NOT NULL,
    credit NUMERIC(3, 1) CHECK (credit > 0),
    dept_id INTEGER NOT NULL REFERENCES department(dept_id)
);
\end{verbatim}

\screenshotplaceholder{6cm}{Screenshot Placeholder: Table creation statements executed successfully}

\subsection*{4.5 Creating a Table with Composite Primary Key}
A composite primary key is formed using two or more columns. It is useful for junction tables.

\begin{verbatim}
CREATE TABLE enrollment (
    student_id INTEGER REFERENCES student(student_id),
    course_id INTEGER REFERENCES course(course_id),
    enrollment_date DATE DEFAULT CURRENT_DATE,
    grade VARCHAR(2),
    PRIMARY KEY (student_id, course_id)
);
\end{verbatim}

\subsection*{4.6 Viewing Table Structure}
In psql, the \texttt{\textbackslash d} command can be used to view the structure of a table.

\begin{verbatim}
\d department
\d student
\d course
\d enrollment
\end{verbatim}

\screenshotplaceholder{6cm}{Screenshot Placeholder: Table descriptions showing columns and constraints}

\subsection*{4.7 Modifying Tables with ALTER}
Use \texttt{ALTER TABLE} to modify an existing table structure.

\noindent Add a new column:
\begin{verbatim}
ALTER TABLE student
ADD COLUMN phone VARCHAR(20);
\end{verbatim}

\noindent Rename a column:
\begin{verbatim}
ALTER TABLE student
RENAME COLUMN phone TO contact_number;
\end{verbatim}

\noindent Change a column data type:
\begin{verbatim}
ALTER TABLE student
ALTER COLUMN contact_number TYPE VARCHAR(30);
\end{verbatim}

\noindent Add a constraint:
\begin{verbatim}
ALTER TABLE enrollment
ADD CONSTRAINT grade_check
CHECK (grade IN ('A+', 'A', 'B+', 'B', 'C+', 'C', 'D', 'F'));
\end{verbatim}

\screenshotplaceholder{6cm}{Screenshot Placeholder: ALTER TABLE statements and updated table structure}

\subsection*{4.8 Renaming and Dropping Database Objects}
Rename a table using \texttt{ALTER TABLE}.

\begin{verbatim}
ALTER TABLE course
RENAME TO offered_course;
\end{verbatim}

Drop a table only when it is no longer needed. The \texttt{IF EXISTS} clause prevents an error if the
table does not exist.

\begin{verbatim}
DROP TABLE IF EXISTS temporary_table;
\end{verbatim}

\subsection*{4.9 Truncating a Table}
The \texttt{TRUNCATE} statement removes all rows from a table while keeping the table structure.

\begin{verbatim}
TRUNCATE TABLE enrollment;
\end{verbatim}

If another table references the table, use \texttt{CASCADE} carefully.

\begin{verbatim}
TRUNCATE TABLE enrollment CASCADE;
\end{verbatim}

\screenshotplaceholder{5cm}{Screenshot Placeholder: TRUNCATE or DROP command execution}

% -------------------------------------------------------------
% 5. Observations
% -------------------------------------------------------------
\section*{5. Observations}
\begin{itemize}
    \item A new PostgreSQL database was created and selected successfully.
    \item A schema was created to organize database objects.
    \item Tables were created with appropriate data types and constraints.
    \item Primary key and foreign key constraints were used to maintain entity integrity and referential integrity.
    \item Existing table structures were modified using \texttt{ALTER TABLE}.
    \item Rows were removed using \texttt{TRUNCATE} without deleting the table definition.
\end{itemize}

% -------------------------------------------------------------
% 6. Lab Task
% -------------------------------------------------------------
\section*{6. Lab Task}

\subsection*{Task A: Library Management Database}
Create a PostgreSQL database named \texttt{library\_lab}. Then create the following tables using DDL
statements:

\begin{enumerate}
    \item \texttt{publisher}: stores publisher information.
    \item \texttt{book}: stores book information and references \texttt{publisher}.
    \item \texttt{member}: stores library member information.
    \item \texttt{borrow}: stores borrowing records and references \texttt{book} and \texttt{member}.
\end{enumerate}

\noindent Requirements:
\begin{itemize}
    \item Use suitable data types for all columns.
    \item Define primary keys for all tables.
    \item Define foreign keys where necessary.
    \item Add at least one \texttt{UNIQUE}, one \texttt{NOT NULL}, one \texttt{CHECK}, and one \texttt{DEFAULT} constraint.
    \item Use \texttt{ALTER TABLE} to add one new column to any table.
    \item Capture screenshots of the SQL commands and table structures.
\end{itemize}

\screenshotplaceholder{7cm}{Screenshot Placeholder: Completed Task A DDL statements and output}

\subsection*{Task B: Hospital Management Database}
Write PostgreSQL DDL statements for a hospital management database containing the following tables:
\texttt{department}, \texttt{doctor}, \texttt{patient}, and \texttt{appointment}.

\noindent Requirements:
\begin{itemize}
    \item Each doctor must belong to a department.
    \item Each appointment must reference one doctor and one patient.
    \item Appointment dates should have a default value of the current date.
    \item Consultation fees must be greater than or equal to zero.
    \item Add suitable primary key, foreign key, unique, not null, default, and check constraints.
    \item Include screenshots of the commands and the final table descriptions.
\end{itemize}

\screenshotplaceholder{7cm}{Screenshot Placeholder: Completed Task B DDL statements and output}

% -------------------------------------------------------------
% 7. Conclusion
% -------------------------------------------------------------
\section*{7. Conclusion}
This experiment introduced PostgreSQL Data Definition Language statements for creating and managing
database structures. Students practiced creating databases, schemas, and tables; defining data types
and constraints; modifying table definitions; and removing or clearing database objects. These DDL
skills provide the foundation for building relational database schemas before inserting and querying
data in later experiments.

\end{document}
